\par\noindent\textbf{(1)} 球表面自由能为 \(4\pi\sigma r^2\)，体相转换释放自由能为 \((4\pi/3)\Delta g\,r^3\)。总自由能变化 \(\Delta G(r)=4\pi\sigma r^2-\frac{4\pi}{3}\Delta g\,r^3\)，二次和三次项的不同幂次产生势垒。
\par\noindent\textbf{(2)} 求导 \(d\Delta G/dr=8\pi\sigma r-4\pi\Delta g\,r^2\)，非零驻点 \(r_c=2\sigma/\Delta g\)。代入得 \(\Delta G_c=\Delta G(r_c)=16\pi\sigma^3/(3\Delta g^2)\)，二阶导数在 \(r_c\) 为负，故其为自由能最大值而不是稳定核。
\par\noindent\textbf{(3)} 临界半径 \(r_c\propto\Delta g^{-1}\)，势垒 \(\Delta G_c\propto\Delta g^{-2}\)。例如把过饱和驱动力增为两倍，临界半径减半、势垒减为四分之一；形核率中指数因子 \(\exp[-\Delta G_c/(k_BT)]\) 因此会对小的热力学驱动变化极其敏感。
\par\noindent\textbf{(4)} 若 \(r<r_c\)，增长方向增加自由能，核倾向缩小；若 \(r>r_c\)，增长方向降低自由能，核倾向长大。真实形核率一般写 \(J=J_0e^{-\Delta G_c/(k_BT)}\)，前因子 \(J_0\) 来自扩散、附着与Zeldovich因子，并非只由静态势垒决定。
